\begin{tabular}
{p{0.4\textwidth} p{0.55\textwidth}}
     \begin{center}
    \fontsize{11}{15}\selectfont ĐẠI HỌC QUỐC GIA TP.HCM\\
         \fontsize{11}{15}\selectfont\textbf{TRƯỜNG ĐẠI HỌC BÁCH KHOA}\\
         -----------------------------
     \end{center}& \begin{center}
         \fontsize{11}{15}\selectfont\textbf{CỘNG HOÀ XÃ HỘI CHỦ NGHĨA VIỆT NAM}\\
         \fontsize{11}{15}\selectfont\textbf{Độc lập - Tự do - Hạnh phúc}\\
         -----------------------------
     \end{center}  
\end{tabular}

\begin{center}
    \fontsize{16}{20}\selectfont
    \textbf{NHIỆM VỤ THỰC TẬP 2} \\

\end{center}
\begingroup
\renewcommand{\arraystretch}{1.5}
\begin{table}[h]
    \centering
    \begin{tabular}{p{0.5\textwidth}p{0.4\textwidth}}
        Họ tên học viên: Văn Xuân Tỷ&MSHV: 2470113  \\
        Ngày, tháng, năm sinh: 05/04/1995 &Nơi sinh: Gia Lai  \\
        Chuyên ngành: Khoa học máy tính &Mã số: 8480101
    \end{tabular}
\end{table}
\endgroup
\subsubsection*{I. TÊN ĐỀ TÀI:}
\onehalfspacing
Phát triển phương pháp lấy mẫu đồ thị cho hệ thống gợi ý quy mô lớn sử dụng mạng nơ-ron đồ thị (Development of Graph Sampling Methods for Large-Scale Recommender Systems Using Graph Neural Networks)

\subsubsection*{II. NHIỆM VỤ VÀ NỘI DUNG}

\begin{itemize}
    \item Hệ thống hóa các nhóm phương pháp lấy mẫu để mở rộng quy mô GNN: lấy mẫu theo đỉnh, theo tầng, theo đồ thị con, và hướng dùng embedding lịch sử.
    \item Trình bày chi tiết định nghĩa, kiến trúc và cơ sở lý thuyết của phương pháp GRAPES.
    \item Tái hiện và phân tích thiết lập thực nghiệm cùng các kết quả chính trong bài báo gốc giới thiệu GRAPES.
    \item Đưa ra nhận định cá nhân về ưu, nhược điểm của phương pháp và đề xuất các hướng cải tiến, gắn với định hướng ứng dụng cho hệ thống gợi ý.\end{itemize}
\subsubsection*{III. NGÀY GIAO NHIỆM VỤ:} 
21/01/2026

\subsubsection*{IV. NGÀY HOÀN THÀNH NHIỆM VỤ:}
12/05/2026

\subsubsection*{V. CÁN BỘ HƯỚNG DẪN:}
TS. Lê Thanh Vân

\vspace{1em}
\begin{tabular}
{p{0.4\textwidth} p{0.5\textwidth}}
\vspace{1.25em}
     \begin{center}
         \textbf{CÁN BỘ HƯỚNG DẪN}
     \end{center}& \begin{center}
     Tp. HCM, ngày 12 tháng 05 năm 2026
     \\
         \textbf{HỘI ĐỒNG NGÀNH}
     \end{center}   \\
\end{tabular}
\vspace{5em}
\begin{center}
    \textbf{TRƯỞNG KHOA KHOA HỌC VÀ KỸ THUẬT MÁY TÍNH}
\end{center}
